\chapter{引言}\label{chap:intro}

% ─────────────────────────────────────────────────────────────────────────────
\section{研究背景与意义}\label{sec:background}

期权是金融衍生品市场的核心工具之一。根据国际清算银行（BIS）统计，全球场外衍生品市场名义本金规模超过 600 万亿美元，其中期权类产品占据重要份额。期权定价的准确性直接影响风险对冲、投资组合管理和监管资本计算，是金融工程的核心问题。

自 Black、Scholes 与 Merton 于 1973 年提出 BSM 模型\citep{black1973pricing,merton1973theory}以来，期权定价理论经历了从常数波动率到局部波动率（CEV 模型\citep{cox1976valuation}）、再到随机波动率（Heston 模型\citep{heston1993closed}）的演进。这三大模型分别对应不同的市场假设，在实践中各有适用场景：BSM 适用于波动率相对平稳的短期定价；CEV 能够捕捉股票市场中波动率随股价下跌而上升的杠杆效应；Heston 则通过引入随机波动率过程，更准确地描述波动率微笑（volatility smile）和厚尾收益率分布。

然而，传统数值方法（有限差分、有限元）在求解高维 PDE 时面临"维数灾难"：网格点数随维度指数增长，Heston 这样的二维模型已需要大量计算资源。此外，传统工具要求用户手动指定模型类型和参数，对非专业用户门槛较高——选择合适的定价模型本身就需要量化金融背景，普通用户难以判断在当前市场条件下应使用 BSM、CEV 还是 Heston。

物理信息神经网络（Physics-Informed Neural Networks，PINN）\citep{raissi2019physics}提供了一种新的思路：将 PDE 残差直接嵌入神经网络的损失函数，无需网格剖分，训练完成后可对任意输入即时推断（单次推断约 2 ms，远低于有限差分的约 500 ms）。近年来，PINN 已被应用于 BSM\citep{dhiman2023pinn}、CEV\citep{bae2024localvol}、Heston\citep{hainaut2024heston,guo2024icpinn}等期权定价模型，展现出良好的精度和效率。

与此同时，大语言模型（Large Language Models，LLM）在自然语言理解和结构化信息提取方面取得了突破性进展。将 LLM 与科学计算工具结合，使用户能够通过自然语言描述问题并获得计算结果，是人机交互的重要方向\citep{liu2025pdeagent}。在期权定价场景中，LLM 的核心价值在于：将"选择合适的定价模型"这一需要专业知识的决策自动化，使非专业用户只需描述期权需求，系统自动完成模型选择和参数提取。

% ─────────────────────────────────────────────────────────────────────────────
\section{三大模型的适用场景与 LLM 路由的必要性}\label{sec:motivation}

三大期权定价模型的应用场景是相互独立的，这正是需要智能路由层的根本原因。

\textbf{BSM 的局限}：BSM 假设波动率为常数，但真实市场中不同行权价的隐含波动率并不相同——这一现象称为波动率微笑（smile）或偏斜（skew）。用 BSM 对深度虚值期权定价会系统性低估其价格，在对冲时产生 Delta 误差。因此，仅支持欧式看涨期权时，BSM 在波动率平稳的市场（如短期国债期权）是足够的，但在股票期权市场中往往不够准确。

\textbf{CEV 的适用场景}：当市场数据显示隐含波动率随行权价降低而升高（负偏斜，常见于股票市场），CEV 模型（$\beta < 1$）能够更好地拟合这一特征。CEV 是 BSM 的自然推广，计算复杂度与 BSM 相当，适合需要捕捉杠杆效应但不需要完整随机波动率建模的场景。

\textbf{Heston 的适用场景}：当市场数据显示波动率本身具有随机性（如 VIX 指数的均值回归特征），且需要同时拟合不同到期日的波动率曲面时，Heston 模型是标准选择。Heston 能够产生厚尾收益率分布，对长期期权和波动率衍生品的定价更为准确，但参数校准需要更多数据。

\textbf{LLM 路由的价值}：上述模型选择逻辑对量化分析师是常识，但对普通用户是门槛。LLM 路由层将这一决策自动化：通过解析用户的自然语言描述（如"股票波动率最近变化很大"→ Heston，"标准欧式期权"→ BSM），自动选择合适的模型并提取参数，使期权定价工具对非专业用户可用。

% ─────────────────────────────────────────────────────────────────────────────
\section{期权类型的扩展性}\label{sec:option_types}

本文当前实现仅支持欧式看涨期权，但框架具有良好的扩展性：

\textbf{欧式看跌期权}：通过 put-call parity 关系
\begin{equation}
  P(S, t) = C(S, t) - S + Ke^{-r(T-t)}
\end{equation}
可直接从看涨价格推导看跌价格，无需额外训练 PINN。

\textbf{美式期权}：美式期权允许持有者在到期日前任意时刻提前行权，其定价问题是一个自由边界问题（free boundary problem），期权价值满足变分不等式：
\begin{equation}
  \min\left(\mathcal{F}[V],\; V - \phi(S)\right) = 0,
\end{equation}
其中 $\phi(S) = \max(S-K, 0)$ 为立即行权价值。PINN 可通过在损失函数中加入互补条件惩罚项处理这类约束，Louskos\citep{louskos2021american} 已验证了这一方法的可行性。

% ─────────────────────────────────────────────────────────────────────────────
\section{研究目标}\label{sec:objectives}

本文旨在构建一个基于 PINN 与 LLM 的统一期权定价框架，实现以下目标：

\begin{enumerate}
  \item 在统一的 PINN 框架下，分别实现 BSM、CEV、Heston 三大期权定价模型的神经网络求解器；
  \item 对 Heston 模型采用 ICPINN 硬约束技术\citep{guo2024icpinn}，确保终值条件精确满足；
  \item 引入 LLM 作为智能路由层，使用户可通过自然语言描述期权需求，由系统自动完成模型选择和参数提取；
  \item 通过与解析解、有限差分参考解的对比，定量评估各 PINN 模型的精度；
  \item 通过参数校准实验，验证框架在真实市场参数下的适用性。
\end{enumerate}

% ─────────────────────────────────────────────────────────────────────────────
\section{本文贡献}\label{sec:contributions}

本文的主要贡献如下：

\begin{itemize}
  \item \textbf{统一框架}：将 BSM、CEV、Heston 三大模型纳入同一 PINN 框架，共用损失函数结构，便于扩展到其他期权定价 PDE。
  \item \textbf{硬约束实现}：对 Heston 模型实现了 ICPINN 输出变换，在二维随机波动率场景下精确满足终值条件，无需在损失函数中加入终值约束项。
  \item \textbf{LLM 路由集成}：将 DeepSeek LLM 与 PINN 求解器结合，通过结构化 JSON 输出实现可靠的参数提取和模型选择，支持中英文自然语言输入，降低期权定价工具的使用门槛。
  \item \textbf{参数校准验证}：基于 SPY 期权链数据对 BSM 和 Heston 模型进行参数校准，验证了框架在真实市场条件下的适用性。
\end{itemize}

% ─────────────────────────────────────────────────────────────────────────────
\section{论文结构}\label{sec:structure}

本文其余部分组织如下：第二章回顾期权定价模型、PINN 方法和 LLM 辅助科学计算的相关工作；第三章介绍本文提出的统一 PINN 框架和 LLM 路由层的设计；第四章报告三大模型的实验结果及参数校准实验；第五章总结全文并讨论局限性与未来工作。

